\hebrewsection{Advanced Topics in Signal Processing}

\subsection{Multiple-Input Multiple-Output (MIMO)}
MIMO systems use multiple antennas at both the transmitter and receiver to improve spectral efficiency or reliability.

\subsubsection{Spatial Multiplexing}
Spatial multiplexing increases the data rate by transmitting independent data streams over different spatial channels. The capacity of a MIMO system with $N_t$ transmit and $N_r$ receive antennas grows linearly with $\min(N_t, N_r)$.

\subsubsection{Beamforming}
Beamforming uses constructive and destructive interference to focus the transmitted energy in a specific direction, increasing the SNR at the receiver and reducing interference for other users.

\subsection{Orthogonal Frequency Division Multiplexing (OFDM)}
OFDM is the backbone of modern high-speed communication (4G, 5G, Wi-Fi). It divides a high-rate data stream into many low-rate streams transmitted on orthogonal subcarriers.

\subsubsection{Cyclic Prefix (CP)}
The CP is a copy of the end of an OFDM symbol prepended to the beginning. It converts linear convolution with the channel into circular convolution, simplifying frequency-domain equalization and eliminates ISI.

\subsection{Spread Spectrum Communications}
Spread spectrum techniques distribute the signal energy over a wide bandwidth.

\subsubsection{Direct Sequence Spread Spectrum (DSSS)}
In DSSS, each bit is multiplied by a high-rate \textbf{pseudo-noise (PN) sequence}. This spreads the signal's bandwidth and provides resistance to narrow-band interference and jamming.

\subsubsection{Frequency Hopping Spread Spectrum (FHSS)}
FHSS changes the carrier frequency rapidly according to a pre-defined sequence.

\subsection{Multi-User Detection (MUD)}
In multi-user systems like CDMA, the receiver can improve performance by jointly detecting all users' signals, effectively cancelling the interference they cause to each other.

\begin{pythonbox}{Simple OFDM Symbol Generation}
\begin{verbatim}
import numpy as np

def generate_ofdm_symbol(data_symbols, n_fft, cp_len):
    # IFFT
    time_domain = np.fft.ifft(data_symbols, n_fft)
    # Add Cyclic Prefix
    ofdm_symbol = np.concatenate([time_domain[-cp_len:], time_domain])
    return ofdm_symbol

# Example
n_fft = 64
cp_len = 16
data = np.random.choice([1+1j, 1-1j, -1+1j, -1-1j], n_fft)
symbol = generate_ofdm_symbol(data, n_fft, cp_len)
\end{verbatim}
\end{pythonbox}

\subsection{The Ambiguity Function}
The ambiguity function $|\chi(\tau, f_d)|$ represents the output of the matched filter in a radar system when the received signal has a delay $\tau$ and a Doppler shift $f_d$. It is a critical tool for waveform design in radar engineering.
